% BOOK-FILE-ID: f-c645fd39fc8a
\subsubsection{整数部分を確定する}\label{節：整数部分を確定する}
ガウス記号があると,そのままでは式変形ができない.
まず考えたいのは,ガウス記号の中身の整数部分を確定させて,記号そのものを外してしまうことである.
整数$k$と実数$x$に対しては
\[
  \lfloor k+x\rfloor=k+\lfloor x\rfloor
\]
が成り立つので,例えば$a_n$が整数なら$\lfloor a_n+\sqrt3\rfloor=a_n+1$と外せる.
中身が平方根なら,隣り合う平方数で挟んでやればよい.
いずれにせよ,先に各項が整数であることを確かめておく必要がある.

\begin{bookpoint}
  \textbf{POINT}\par
  中身を$k\le X<k+1$の形に挟んでガウス記号を外し,既知の漸化式に帰着させる.
\end{bookpoint}

\noindent \bookblocklink{example-gauss-role-1}{problem}{answer}{【例題】}\par\nobreak
\begin{enumerate}[label=(\arabic*),parsep=3pt,beginpenalty=10000,format=\bookitemlink{example-gauss-role-1}{problem}{answer}]
  \item \bookterm{sequence}{数列}を次のように定める.\bookterm{general}{一般項}を求めよ.
    \BookMathLayout{            \[
        a_1=1,\qquad
        a_{n+1}=\left\lfloor\sqrt{a_n^2+4a_n+5}\right\rfloor
        \quad(n\ge1).
      \]}{            \[
        \begin{aligned}
          a_1&=1,\\
          a_{n+1}&=\left\lfloor\sqrt{a_n^2+4a_n+5}\right\rfloor
          \\              &\quad(n\ge1).
        \end{aligned}
      \]}
\end{enumerate}

\noindent \bookblocklink{example-gauss-role-1}{hint}{problem}{【方針】}\par\nobreak
\begin{enumerate}[label=(\arabic*),parsep=3pt,beginpenalty=10000,format=\bookitemlink{example-gauss-role-1}{hint}{problem}]
  \item 整数$a_n$に対して,平方根の中身を$(a_n+2)^2$と$(a_n+3)^2$で挟みます.
    まず各項が非負整数であることを確かめましょう.
\end{enumerate}

\noindent \bookblocklink{example-gauss-role-1}{answer}{problem}{【解答】}\par\nobreak
\begin{enumerate}[label=(\arabic*),parsep=3pt,beginpenalty=10000,format=\bookitemlink{example-gauss-role-1}{answer}{problem}]
  \item \[
      a_n=2n-1\quad(n\ge1).
    \]
\end{enumerate}

\noindent \bookblocklink{example-gauss-role-1}{solution}{problem}{【解法】}\par\nobreak
\begin{enumerate}[label=(\arabic*),parsep=3pt,beginpenalty=10000,format=\bookitemlink{example-gauss-role-1}{solution}{problem}]
  \item 初項と漸化式から,帰納的に各項は非負整数となる.
    非負整数$x$について,
    \[
      (x+2)^2\le x^2+4x+5<(x+3)^2
    \]
    である(左の差は1,右の差は$2x+4>0$).平方根を取ると,
    \[
      x+2\le\sqrt{x^2+4x+5}<x+3
    \]
    より$a_{n+1}=a_n+2$.\,\sectionref{節：等差型}に帰着し,\,$a_n=2n-1$を得る.
\end{enumerate}
